% ================================================================
%  conteudo/semana03.tex
% ================================================================

\section{Semana 3 --- 30 de março de 2026 a 2 de abril de 2026}

\begin{table}[H]
\centering
\caption{Registro de atividades da Semana 3.}
\begin{tabular}{llll}
\toprule
\textbf{Data} & \textbf{Dia} & \textbf{Horas} & \textbf{Descrição} \\
\midrule
30/03/2026 & Segunda-feira & 2,0 h & Atividades no CIA \\
31/03/2026 & Terça-feira & 2,0 h & Atividades no CIA \\
01/04/2026 & Quarta-feira & 4,0 h & Atividades no CIA \\
02/04/2026 & Quinta-feira & 4,0 h & Atividades no CIA \\
03/04/2026 & Sexta-feira & -- & Sexta-feira Santa (sem atividade) \\
\midrule
\multicolumn{2}{l}{\textbf{Total da semana}} & \textbf{12,0 h} & \\
\multicolumn{2}{l}{\textbf{Total acumulado}} & \textbf{47,0 h} & \\
\bottomrule
\end{tabular}
\end{table}

\subsubsection*{Atividades realizadas}

A terceira semana seguiu a rotina consolidada de análise e preparo de
materiais, com destaque para duas atividades específicas relacionadas à
infraestrutura criogênica do laboratório:

\begin{itemize}
    \item Continuidade das rotinas de análise e preparo de materiais para
          IRMS, TGA, RMN, BET e DSC;
    \item Manutenção e reabastecimento dos gases dos equipamentos;
    \item \textbf{Carregamento de gás Hélio} nos tanques criogênicos do
          laboratório, procedimento essencial para a manutenção do campo
          magnético supercondutor do espectrômetro de RMN;
    \item \textbf{Primeiros testes do \textit{script}} de monitoramento de
          tanques criogênicos (He e $\text{N}_2$) desenvolvido na semana
          anterior, verificando a leitura correta dos níveis e a
          estabilidade do programa instalado no computador do RMN.
\end{itemize}

\subsubsection*{Observações}

O carregamento de hélio é uma operação crítica e periódica no CIA,
exigindo atenção ao manuseio seguro do criogênico (temperatura de
$-268{,}9\,^{\circ}$C). Os primeiros testes do \textit{script} de
monitoramento permitiram validar a lógica de leitura e identificar
pontos de ajuste no código, iniciando o ciclo de testes e refinamento
da ferramenta computacional desenvolvida.
